\section*{Bab \ref{ch:g8:winds-global-circulation} --- Angin dan Sirkulasi Global}

\begin{solution}{exo:g8:winds-global-circulation:1}
(a) \omterm{def:g8:winds-global-circulation:global-winds}{Angin pasat} timur laut, dari timur laut. (b) \omterm{def:g8:winds-global-circulation:global-winds}{Angin baratan}, dari barat
laut. (c) \omterm{def:g8:winds-global-circulation:global-winds}{Angin timuran kutub}, dari timur (dari timur laut). (d) \omterm{def:g8:winds-global-circulation:global-winds}{Angin
pasat} tenggara, dari tenggara.
\end{solution}

\begin{solution}{exo:g8:winds-global-circulation:2}
(a) Udara hangat naik dan udara dingin turun. (b) Rendah. (c) Dari tekanan
tinggi ke tekanan rendah. (d) Ke kanan.
\end{solution}

\begin{solution}{exo:g8:winds-global-circulation:3}
(a) Salah: \omterm{def:g8:winds-global-circulation:breezes}{angin laut} bertiup dari \omterm{def:g2:water-on-earth:ocean}{laut} ke darat (pada siang hari).
(b) Benar. (c) Salah: \omterm{def:g8:winds-global-circulation:global-winds}{angin pasat} bertiup dari timur (dari timur laut atau
dari tenggara). (d) Benar.
\end{solution}

\begin{solution}{exo:g8:winds-global-circulation:4}
Gambarnya menunjukkan sebuah lingkaran: di sepanjang tanah, panah dari \omterm{def:g2:water-on-earth:ocean}{laut}
ke darat; di atas daratan, panah yang naik; di \omterm{def:g4:world-physical-map:altitude}{ketinggian}, panah dari darat
kembali ke atas \omterm{def:g2:water-on-earth:ocean}{laut}; di atas \omterm{def:g2:water-on-earth:ocean}{laut}, panah yang turun. T ditandai di tanah
di bawah udara yang turun (di atas \omterm{def:g2:water-on-earth:ocean}{laut}), R di bawah udara yang naik (di
atas daratan).
\end{solution}

\begin{solution}{exo:g8:winds-global-circulation:5}
(a) Tanah 1.41 kali lebih luas. (b) 50~\%. (c) 26~\%. (d) Di dekat \omterm{def:g3:continents-and-oceans:poles}{kutub},
sinar matahari tiba dengan miring dan tersebar pada wilayah yang luas,
sehingga setiap meter persegi hanya menerima sedikit energi.
\end{solution}

\begin{solution}{exo:g8:winds-global-circulation:6}
(a) \omterm{def:g8:winds-global-circulation:global-winds}{Angin timuran kutub} (\ang{70} sampai \ang{60}\,LU), \omterm{def:g8:winds-global-circulation:global-winds}{angin baratan}
(\ang{60} sampai \ang{30}\,LU), \omterm{def:g8:winds-global-circulation:global-winds}{angin pasat} timur laut (\ang{30}\,LU
sampai \omterm{def:g3:continents-and-oceans:equator}{khatulistiwa}), \omterm{def:g8:winds-global-circulation:global-winds}{angin pasat} tenggara (\omterm{def:g3:continents-and-oceans:equator}{khatulistiwa} sampai
\ang{30}\,LS), \omterm{def:g8:winds-global-circulation:global-winds}{angin baratan} (\ang{30} sampai \ang{60}\,LS), \omterm{def:g8:winds-global-circulation:global-winds}{angin timuran
kutub} (\ang{60} sampai \ang{70}\,LS). (b) Di dekat \omterm{def:g3:continents-and-oceans:equator}{khatulistiwa}, di daerah
doldrum. (c) Di dekat \ang{30}\,LU dan \ang{30}\,LS, di bawah udara yang
turun.
\end{solution}

\begin{solution}{exo:g8:winds-global-circulation:7}
(a) Daratan ($31 > 22$): angin bertiup dari \omterm{def:g2:water-on-earth:ocean}{laut} ke darat, yaitu \omterm{def:g8:winds-global-circulation:breezes}{angin
laut}. (b) \omterm{def:g2:water-on-earth:ocean}{Laut} ($21 > 17$): angin bertiup dari darat ke \omterm{def:g2:water-on-earth:ocean}{laut}, yaitu \omterm{def:g8:winds-global-circulation:breezes}{angin
darat}. (c) Daratan: $31 - 17 = 14$ derajat; \omterm{def:g2:water-on-earth:ocean}{laut}: $22 -
21 = 1$ derajat. (d) Sinar Matahari hanya memanaskan selapis tipis tanah,
yang cepat menjadi panas dan cepat menjadi dingin; \omterm{def:g2:water-on-earth:ocean}{laut} menyebarkan
panasnya ke lapisan air yang tebal, sehingga suhunya hampir tidak berubah
antara siang dan malam.
\end{solution}

\begin{solution}{exo:g8:winds-global-circulation:8}
(a) Dari utara. (b) Ke kanan. (c) Dari timur laut (itulah \omterm{def:g8:winds-global-circulation:global-winds}{angin pasat}
timur laut). (d) Tanpa rotasi, udara itu akan bertiup dari selatan; rotasi
membelokkannya ke kiri; udara itu bertiup dari tenggara.
\end{solution}

\begin{solution}{exo:g8:winds-global-circulation:9}
Di dekat \omterm{def:g3:continents-and-oceans:equator}{khatulistiwa}, tanah dan \omterm{def:g2:water-on-earth:ocean}{laut} \emph{dipanaskan} dengan kuat oleh
Matahari; udara yang hangat dan lembap \emph{naik}; selagi naik, udara itu
\emph{mendingin}, dan uap airnya berubah menjadi \emph{\omterm{def:g4:clouds-rain-snow:cloud}{awan}}, yang
menurunkan \omterm{def:g4:clouds-rain-snow:precipitation}{hujan} deras dan badai petir. Di dekat \ang{25}\,LU, udara yang
naik di \omterm{def:g3:continents-and-oceans:equator}{khatulistiwa} \emph{turun}; di bawahnya, tekanan di dekat tanah
\emph{tinggi}; udara yang turun itu menghangat, \omterm{def:g4:clouds-rain-snow:cloud}{awannya} mengering, dan
wilayah itu tetap \emph{kering}: sebuah \omterm{def:g5:weather-and-climate:desert}{gurun}.
\end{solution}

\begin{solution}{exo:g8:winds-global-circulation:10}
(a) \omterm{def:g8:winds-global-circulation:global-winds}{Angin baratan}. Tidak: angin itu bertiup dari barat, berlawanan dengan
kapal yang berlayar ke barat. (b) Untuk mencapai sabuk \omterm{def:g8:winds-global-circulation:global-winds}{angin pasat}, di
selatan kira-kira \ang{30}\,LU. (c) \omterm{def:g8:winds-global-circulation:global-winds}{Angin pasat} timur laut, antara
\omterm{def:g3:continents-and-oceans:equator}{khatulistiwa} dan kira-kira \ang{30}\,LU. (d) \omterm{def:g8:winds-global-circulation:global-winds}{Angin baratan}, di utara
kira-kira \ang{30}\,LU (di dekat \ang{40}\,LU, \omterm{def:g6:latitude-longitude:latitude}{lintang} pelabuhan asalnya).
\end{solution}

\begin{solution}{exo:g8:winds-global-circulation:11}
(a) $1 \div 0.5 = 2$: tabelnya memberikan 2. (b) $1 \div 0.866 \approx
1.1547$, yaitu \textbf{1.15}. (c) Pada \ang{30}: $1000 \div 1.15 \approx
\qty{870}{W}$ pada setiap meter persegi (tepatnya $1000 \times 0.866 =
\qty{866}{W}$). Pada \ang{60}: $1000 \div 2 = \qty{500}{W}$.
\end{solution}

\begin{solution}{exo:g8:winds-global-circulation:12}
(a) Tiga lingkaran: udara naik di \omterm{def:g3:continents-and-oceans:equator}{khatulistiwa} dan turun di dekat
\ang{30}; naik di dekat \ang{60} dan turun di dekat \ang{30}; naik di dekat
\ang{60} dan turun di \omterm{def:g3:continents-and-oceans:poles}{kutub}. (b) R di \omterm{def:g3:continents-and-oceans:equator}{khatulistiwa} dan di dekat \ang{60};
T di dekat \ang{30} dan di \omterm{def:g3:continents-and-oceans:poles}{kutub}. (c) Dari \ang{30} ke \omterm{def:g3:continents-and-oceans:equator}{khatulistiwa}, udara
di dekat permukaan bergerak ke selatan dan dibelokkan ke kanan: \omterm{def:g8:winds-global-circulation:global-winds}{angin pasat}
timur laut. Dari \ang{30} ke \ang{60}, udara itu bergerak ke utara dan
dibelokkan ke kanan: udara itu datang dari barat daya. Dari \omterm{def:g3:continents-and-oceans:poles}{kutub} ke
\ang{60}, udara itu bergerak ke selatan, dibelokkan ke kanan: \omterm{def:g8:winds-global-circulation:global-winds}{angin timuran
kutub}. (d) Di sel yang tengah, udara di dekat permukaan mengalir ke arah
\omterm{def:g3:continents-and-oceans:poles}{kutub}, yaitu ke utara; setelah dibelokkan ke kanan, udara itu menuju ke
timur laut: itulah \omterm{def:g8:winds-global-circulation:global-winds}{angin baratan}.
\end{solution}

\begin{solution}{pb:g8:winds-global-circulation:1}
\textbf{1.} $59.6 - 23.5 = 36.1$ derajat \omterm{def:g6:latitude-longitude:longitude}{bujur}.

\textbf{2.} $36.1 \times 108 \approx 3899$: kira-kira \qty{3900}{km}.

\textbf{3.} Ya: $3900 - 3896 = 4$~km, selisih yang sangat kecil.

\textbf{4.} Ya: keduanya terletak dekat \ang{14}\,LU, antara \omterm{def:g3:continents-and-oceans:equator}{khatulistiwa}
dan kira-kira \ang{30}\,LU, sabuk \omterm{def:g8:winds-global-circulation:global-winds}{angin pasat} timur laut.

\textbf{5.} Dari timur laut (panah-panahnya mengarah ke barat daya).

\textbf{6.} Ke depan, dan sedikit ke samping: angin itu datang dari
belakang kapal, sedikit dari sebelah kanannya, dan mendorongnya ke arah
barat.

\textbf{7.} Angin itu bagian dari \omterm{def:g8:winds-global-circulation:circulation}{sirkulasi global}: udara terus turun di
dekat \ang{30}\,LU dan naik di dekat \omterm{def:g3:continents-and-oceans:equator}{khatulistiwa}, sehingga udara di dekat
permukaan terus mengalir ke arah \omterm{def:g3:continents-and-oceans:equator}{khatulistiwa}, dan rotasi Bumi selalu
membelokkannya ke kanan.

\textbf{8.} \omterm{def:g8:winds-global-circulation:global-winds}{Angin pasat} akan bertiup berlawanan dengan kapal itu. Nakhoda
harus berlayar ke utara, melewati kira-kira \ang{30}\,LU, untuk menemukan
\omterm{def:g8:winds-global-circulation:global-winds}{angin baratan}.

\textbf{9.} $7 \times 24 = \qty{168}{km}$ sehari.

\textbf{10.} $3 \times 168 = \qty{504}{km}$.

\textbf{11.} $3896 \div 168 \approx 23.2$ hari.

\textbf{12.} Kira-kira \textbf{23 hari}: kurang dari sebulan.
\end{solution}
